\setcounter{chapter}{6}\chapter{Primitives et équations différentielles}\label{ch:7}

{\small\sffamily\color{teal}S52, S53, S55, S58\par}

\begin{retenir}{Objectifs}Calculer des primitives ; résoudre une équation différentielle avec condition initiale ; interpréter une approximation d'Euler.\end{retenir}

\section{Notes de cours}
\subsection{Primitives sur un intervalle}
On admet à ce stade que toute fonction continue sur un intervalle admet des primitives ; le lien avec l'aire sera établi au chapitre 8.
Une primitive de $f$ sur un intervalle $I$ est une fonction dérivable $F$ telle que $F'=f$ sur $I$. Si $F$ est une primitive, toutes les primitives sur $I$ sont $F+C$, $C\in\R$.
\paragraph{Démonstration.} Si $F'=G'=f$, alors $(F-G)'=0$ sur l'intervalle ; $F-G$ est donc constante. Réciproquement, $(F+C)'=f$. L'hypothèse d'intervalle est essentielle : sur un domaine disjoint, les constantes peuvent différer selon les intervalles.
\begin{retenir}{Table de primitives --- à lire en sens inverse des dérivées}
\[\begin{array}{c|c}
f(x)&F(x)\\\hline
x^n\ (n\in\Z,\ n\ne-1)&\frac{x^{n+1}}{n+1}\\
\frac{1}{x}&\ln|x|\quad\text{sur un intervalle évitant }0\\
\frac{1}{\sqrt{x}}&2\sqrt{x}\quad(x>0)\\
e^x&e^x\\
\cos x&\sin x\\
\sin x&-\cos x
\end{array}\]
La constante additive peut toujours être ajoutée. Les deux dernières lignes seront réinvesties après l'étude de la trigonométrie.
\end{retenir}
\subsection{Reconnaître une composition}
Pour une fonction $u$ dérivable, les primitives de $u'u^n$, $u'e^u$, $\frac{u'}{u}$ sont respectivement $\frac{u^{n+1}}{n+1}$ ($n\ne-1$), $e^u$, $\ln|u|$, sur les intervalles où les expressions ont un sens. La linéarité permet de traiter les sommes.
\paragraph{Exemple.} Une primitive de $6x(x^2+1)^2$ est $(x^2+1)^3$ ; une primitive de $e^{3x}$ est $\frac{e^{3x}}{3}$. Pour imposer $F(0)=2$ dans ce dernier cas, prendre $F(x)=\frac{e^{3x}}{3}+\frac{5}{3}$.
\subsection{Équations différentielles}
Une équation différentielle relie une fonction inconnue à ses dérivées. Une solution doit être vérifiée sur un intervalle précisé.
\begin{retenir}{Solutions usuelles}
$y'=ay$ a pour solutions $y(x)=Ce^{ax}$. Pour $a\ne0$, $y'=ay+b$ a pour solutions
\[y(x)=Ce^{ax}-\frac ba.\]
Si $a=0$, l'équation devient $y'=b$ et $y(x)=bx+C$.
Si $g$ est une solution particulière connue de $y'=ay+f(x)$, les solutions sont $y=g+Ce^{ax}$.
\end{retenir}
\paragraph{Démonstration pour $y'=ay$.} Si $y$ est solution, $z(x)=e^{-ax}y(x)$ vérifie $z'=e^{-ax}(y'-ay)=0$ ; $z=C$. Réciproquement, la dérivée de $Ce^{ax}$ est $aCe^{ax}$. Pour $y'=ay+b$, soustraire la solution constante $-\frac{b}{a}$ ramène à l'équation homogène. Pour un second membre variable, faire de même avec $g$.
\paragraph{Exemple.} $y'=-2y+6$ et $y(0)=1$ donnent $y=3+Ce^{-2x}$, puis $C=-2$. La solution $3-2e^{-2x}$ croît vers l'équilibre $3$ lorsque $x\to+\infty$.
\subsection{Méthode d'Euler}
Pour $y'=ay+b$, avec un pas $h>0$, on approche $y(t_k)$ par $Y_k$ selon $t_{k+1}=t_k+h$, $Y_{k+1}=Y_k+h(aY_k+b)$. C'est une approximation utilisant la tangente, pas une égalité sur la solution exacte. Un pas plus petit peut améliorer l'approximation ; contrôler numériquement l'erreur sur une solution connue.


\section{Activité de recherche}\textbf{Refroidissement d'une boisson --- 40 min.}\par Une boisson initialement à $80\,^{\circ}\mathrm C$ est placée dans une pièce maintenue à $20\,^{\circ}\mathrm C$. Pour $t\ge0$, on note $T(t)$ sa température en degrés Celsius, $t$ désignant le temps écoulé en minutes. On admet le modèle $T' (t)=-0{,}1(T(t)-20)$ et $T(0)=80$.
\begin{enumerate}\item Résoudre ce problème initial et donner l'expression de $T(t)$.
\item Étudier le sens de variation de $T$ et déterminer sa limite en $+\infty$. Interpréter les résultats.
\item Déterminer le premier instant où la température vaut $40\,^{\circ}\mathrm C$. Donner la valeur exacte, puis une valeur approchée en minutes.
\item Appliquer deux étapes de la méthode d'Euler avec un pas de $1$ minute. Comparer l'approximation obtenue au bout de deux minutes à la valeur exacte $T(2)$.
\item Indiquer une situation concrète dans laquelle l'hypothèse d'une température ambiante constante ne serait plus adaptée.\end{enumerate}

\section{Exercices progressifs}

\exercice{1}{Entraînement}Déterminer toutes les primitives sur $\R$ de la fonction $x\mapsto3x^2-4x+1$.

\exercice{2}{Entraînement}Déterminer la primitive $F$ de la fonction $x\mapsto\frac{2x}{x^2+4}$ sur $\R$ qui vérifie $F(0)=0$.

\exercice{3}{Entraînement}Déterminer la solution sur $\R$ de l'équation différentielle $y' =3y$ qui vérifie $y(1)=2$.

\exercice{4}{Entraînement}Déterminer la solution sur $\R$ de l'équation différentielle $y' =-y+4$ qui vérifie $y(0)=7$.

\exercice{5}{Entraînement}On considère l'équation différentielle $(E):y' =y-x$ sur $\R$. Vérifier que $g(x)=x+1$ en est une solution, puis déterminer la solution de $(E)$ qui vérifie $y(0)=0$.

\section{Travail dans le manuel ou une annale}\textbf{Barbazo, p. 228, exercices 25 et 30.} En S55 : exercice 30, pour faire apparaître une solution constante. À la maison : exercice 25 en entier, en distinguant une condition sur $f$ d'une condition sur $f'$.

\section{Devoir à la maison}\textbf{DM7 --- Une évolution avec apport variable --- 60 min.}\par On cherche une fonction $y$ définie et dérivable sur $[0;+\infty[$ telle que $y' (t)=-y(t)+2t+1$ et $y(0)=3$.
\begin{enumerate}\item Déterminer deux réels $a$ et $b$ pour que $g(t)=at+b$ soit une solution de l'équation différentielle.
\item En déduire l'expression de la solution qui vérifie $y(0)=3$.
\item Étudier les variations de cette solution et préciser son minimum ainsi que l'instant où il est atteint.
\item Écrire une fonction Python d'argument $N$, entier strictement positif, qui applique la méthode d'Euler sur $[0;2]$ avec le pas $\frac2N$ et renvoie l'approximation de $y(2)$. Comparer avec la valeur exacte pour $N=10$, puis $N=100$.
\item Pourquoi la vérification de l'égalité différentielle en quelques instants seulement ne suffit-elle pas à prouver qu'une fonction est solution ?\end{enumerate}

\ifprof\clearpage\section{Conduite et corrigés}\begin{profbox}{Repères pour le professeur}Distinguer fonction inconnue et variable réelle. Faire vérifier chaque primitive par dérivation. Pour $y'=ay+b$, chercher d'abord la solution constante. Le cas $a=0$ doit rester visible. Les primitives trigonométriques sont données ici comme table de référence, mais leur utilisation évaluée attend la séquence de trigonométrie. Pour Euler, garder les unités et comparer à une solution exacte connue.\end{profbox}

\subsection{Corrigé de l'activité}\begin{enumerate}
\item L'équation est $T'=-0{,}1T+2$. Une solution constante est $20$ ; toutes les solutions sont $T(t)=20+Ce^{-0{,}1t}$. Comme $T(0)=80$, $C=60$. On vérifie $T'(t)=-6e^{-0{,}1t}=-0{,}1(T(t)-20)$ et $T(0)=80$.
\item $T'(t)=-6e^{-0{,}1t}<0$ : la température décroît strictement. Comme $e^{-0{,}1t}\to0$ en $+\infty$, $T(t)\to20\,^{\circ}\mathrm C$, en restant supérieure à $20\,^{\circ}\mathrm C$.
\item $T(t)=40\Longleftrightarrow e^{-0{,}1t}=\frac{1}{3}$. En prenant le logarithme, $-0{,}1t=-\ln3$, soit $t=10\ln3\simeq10{,}9861$ min. La stricte décroissance assure qu'il s'agit de l'unique instant de passage à $40\,^{\circ}\mathrm C$.
\item Pour $h=1$, $Y_{k+1}=Y_k-0{,}1(Y_k-20)$. Ainsi $Y_1=80-6=74$ et $Y_2=74-5{,}4=68{,}6$. La solution exacte donne $T(2)=20+60e^{-0{,}2}\simeq69{,}1238$. Euler sous-estime ici la température d'environ $0{,}5238\,^{\circ}\mathrm C$.
\item Le modèle suppose une température ambiante de $20\,^{\circ}\mathrm C$ et un coefficient d'échange constant de $0{,}1\,\mathrm{min}^{-1}$. Un changement de pièce, un couvercle ou une agitation peut modifier ces hypothèses ; la formule n'est alors plus automatiquement adaptée.
\end{enumerate}

\paragraph{E1.}La linéarité donne $F(x)=x^3-2x^2+x+C$, $C\in\R$. On vérifie $F'(x)=3x^2-4x+1$. Toutes les primitives sur l'intervalle $\R$ diffèrent d'une constante.

\paragraph{E2.}Le dénominateur est positif et sa dérivée vaut $2x$. Les primitives sont $\ln(x^2+4)+C$. La condition en $0$ donne $C=-\ln4$, donc $F(x)=\ln(x^2+4)-\ln4$ ; on vérifie $F(0)=0$ et $F'(x)=\frac{2x}{x^2+4}$.

\paragraph{E3.}Les solutions de $y'=3y$ sont $Ce^{3x}$. La condition $Ce^3=2$ impose $C=2e^{-3}$. Ainsi $y(x)=2e^{3(x-1)}$, dont la dérivée vaut $3y(x)$ et dont la valeur en $1$ est $2$.

\paragraph{E4.}La constante $4$ est une solution particulière. Les solutions sont donc $4+Ce^{-x}$. Comme $4+C=7$, $C=3$. La solution est $y(x)=4+3e^{-x}$ ; sa dérivée $-3e^{-x}$ est bien $-y+4$, et $y(0)=7$.

\paragraph{E5.}$g(x)=x+1$ vérifie $g'=1=g-x$. Si $z=y-g$, alors $z'=z$, donc $y=x+1+Ce^x$. La condition $y(0)=0$ impose $C=-1$. On obtient $y=x+1-e^x$ et l'on vérifie $y'=1-e^x=y-x$.

\subsection{Corrigé du DM7}\begin{enumerate}
\item Si $g(t)=at+b$, alors $g'=a$ et $-g+2t+1=(2-a)t+1-b$. Pour que l'identité soit vraie à tout $t$, il faut $2-a=0$ et $a=1-b$. Donc $a=2$, $b=-1$, et $g(t)=2t-1$.
\item Posons $z=y-g$. En soustrayant les deux équations, $z'=-z$, donc $z=Ce^{-t}$. Ainsi $y(t)=2t-1+Ce^{-t}$. Comme $y(0)=3$, $C=4$. La fonction $y(t)=2t-1+4e^{-t}$ vérifie bien $y'=2-4e^{-t}=-y+2t+1$ et $y(0)=3$.
\item $y'(t)=2-4e^{-t}$. Cette dérivée est négative pour $t<\ln2$, nulle en $\ln2$, positive après. La solution décroît sur $[0,\ln2]$ puis croît sur $[\ln2,+\infty[$. Son minimum est $y(\ln2)=2\ln2-1+2=1+2\ln2$.
\item 
\begin{lstlisting}
def euler_dm7(N):
    h = 2/N
    y = 3.0
    for k in range(N):
        t = k*h
        y = y+h*(-y+2*t+1)
    return y
\end{lstlisting}
Le temps utilisé dans le second membre est celui du début du pas. Le programme donne $Y_{10}\simeq3{,}42950$ et $Y_{100}\simeq3{,}53048$, alors que $y(2)=3+4e^{-2}\simeq3{,}54134$. Les erreurs absolues valent environ $0{,}11184$ et $0{,}01086$.
Pour vérifier ces valeurs sans boucle, posons $Z_k=Y_k-(2kh-1)$. La récurrence d'Euler donne $Z_{k+1}=(1-h)Z_k$ et $Z_0=4$. Donc $Y_N=3+4(1-\frac{2}{N})^N$.
\item Une solution d'équation différentielle doit vérifier une égalité entre fonctions sur tout l'intervalle, ainsi que la condition initiale. Quelques valeurs numériques ne contrôlent qu'un nombre fini de points ; la substitution symbolique de la question 2 est une vérification sur tout le domaine.
\end{enumerate}

\subsection{Corrigé du travail dans le manuel ou une annale}\paragraph{Exercice 30.} Si $y$ est constante égale à $\ell$, l'équation devient $0=-\frac12\ell+3$, donc $\ell=6$. En posant $z=y-6$, on obtient $z'=-\frac12z$ ; toutes les solutions sont $y(x)=6+Ce^{-\frac x2}$, avec $C\in\R$.
\paragraph{Exercice 25.}
\begin{enumerate}
\item $f'=\frac12f$ donne $f(x)=Ce^{\frac x2}$. Comme $f(0)=1$, $C=1$.
\item $5f'=f$ donne $f(x)=Ce^{\frac x5}$. La condition $f(1)=0$ impose $Ce^{\frac15}=0$ ; l'exponentielle n'étant jamais nulle, $C=0$. La solution est la fonction nulle.
\item $f'+f=0$ donne $f(x)=Ce^{-x}$, donc $f'(x)=-Ce^{-x}$. La condition porte sur la dérivée : $f'(1)=2$ équivaut à $-Ce^{-1}=2$, soit $C=-2e$. Ainsi $f(x)=-2e^{1-x}$.
\end{enumerate}

\fi

